\subsection{Disintegration of a measure by a measurable equivalence relation}

Let T be a Polish locally compact space, \(\mu\) a positive measure on T, and R a \(\mu\) -measurable equivalence relation in T. Then, there exist (No.~4, Prop. 2) a Polish locally compact space B and a \(\mu\) -measurable mapping p of T into B, such that the relation \(p(x) = p(y)\) is equivalent to \(R\{x, y\}\) . Every measure \(\nu\) that is a pseudo-image of \(\mu\) under p (No.~2) will be called a quotient measure of \(\mu\) by the relation R; if \(b \mapsto \lambda_b\) is a disintegration of \(\mu\) relative to the measure \(\nu\) , we shall say that \(b \mapsto \lambda_b\) is a disintegration of \(\mu\) by the relation R. By virtue of the properties of p and the \(\lambda_b\) , each of the measures \(\lambda_b\) is concentrated on an equivalence class with respect to R, and if \(b \neq c\) , the measures \(\lambda_b\) and \(\lambda_c\) are concentrated on distinct classes.

The space B, the mapping p and the pseudo-image measure \(\nu\) on B can in general be chosen in infinitely many ways. Nevertheless, the various disintegrations of \(\mu\) by R can all be deduced from one among them, as a consequence of the following theorem:

THEOREM 4. --- Let T be a Polish locally compact space, \(\mu\) a positive measure on T, and R a \(\mu\) -measurable equivalence relation in T. Let B, B' be two Polish locally compact spaces, \(p, p'\) two \(\mu\) -measurable mappings of T into B, B' respectively, such that R\{x,y\} is equivalent to \(p(x) = p(y)\) and to \(p'(x) = p'(y)\) . Let \(\nu, \nu'\) be pseudo-image measures of \(\mu\) under \(p, p'\) respectively; let \(b \mapsto \lambda_b\) , \(b' \mapsto \lambda_{b'}'\) be disintegrations of \(\mu\) relative to \(\nu, \nu'\) respectively.

Under these conditions, there exist in B (resp. B') a set N (resp. N') negligible for \(\nu\) (resp. \(\nu'\) ) and a bijection f of B - N onto B' - N', having the following properties:

\begin{enumerate}
\def\labelenumi{\alph{enumi})}
\item
  The mapping \(f\) (defined almost everywhere in B) is \(\nu\) -measurable and its inverse mapping \(f'\) is \(\nu'\) -measurable; every pseudo-image measure of \(\nu\) (resp. \(\nu'\) ) under \(f\) (resp. \(f'\) ) is equivalent to \(\nu'\) (resp. \(\nu\) ).
\item
  For every \(b \in B - N\) , the measure \(\lambda_{f(b)}^{\prime}\) on T is of the form \(r(b)\lambda_{b}\) where \(r(b) \neq 0\) and r is locally \(\nu\) -integrable.
\end{enumerate}

To establish \(a\) ), we may limit ourselves to the case that \(\nu\) and \(\nu'\) are bounded measures (Ch. V, §5, No.~6, Prop. 11). Let \(\mathrm{N}_0 = \mathrm{B} - p(\mathrm{T})\) , \(\mathrm{N}_0' = \mathrm{B}' - p'(\mathrm{T})\) ; we know that \(\mathrm{N}_0\) (resp. \(\mathrm{N}_0'\) ) is negligible for \(\nu\) (resp. \(\nu'\) ) (No.~2). There exists a bijection \(f\) of \(\mathrm{B} - \mathrm{N}_0\) onto \(\mathrm{B}' - \mathrm{N}_0'\) defined by \(f(p(t)) = p'(t)\) for all \(t \in T\) ; let \(f'\) be the inverse mapping of \(f\) , so that \(f'(p'(t)) = p(t)\) . For every subset \(M\) of \(B\) , the relation « M is \(\nu\) -measurable» is equivalent to « \(\overline{p}^{-1}(M)\) is \(\mu\) -measurable», that is, to « \(p'(f(M))\) is \(\mu\) -measurable», thus finally to « \(f(M)\) is \(\nu'\) -measurable» (Ch. V, §6, No.~2, Cor. of Prop. 3). We thus see that \(f\) (resp. \(f'\) ) transforms every \(\nu\) -measurable (resp. \(\nu'\) -measurable) set into a \(\nu'\) -measurable (resp. \(\nu\) -measurable) set; since B and \(B'\) are metrizable and have countable bases, it follows that \(f\) and \(f'\) are measurable (Ch. IV, §5, No.~5, Th. 4). Moreover, if \(M \subset B\) is \(\nu\) -negligible then \(\overline{p}^{-1}(M) = p'(f(M))\) is \(\mu\) -negligible, therefore \(f(M)\) is \(\nu'\) -negligible (Ch. V, §6, No.~2, Cor. 2 of Prop. 2); similarly, \(f'\) transforms every \(\nu'\) -negligible set into a \(\nu\) -negligible set. Consequently, the image of \(\nu\) under \(f\) (which is defined since \(\nu\) is bounded, which implies that \(f\) is \(\nu\) -proper) is equivalent to \(\nu'\) , and the image of \(\nu'\) under \(f'\) is equivalent to \(\nu\) (Ch. V, §5, No.~6, Prop. 10). It remains to prove \(b\) ). By virtue of Th. 2 of No.~3, we can restrict ourselves to the case that \(\nu' = f(\nu)\) . Since \(\mu = \int \lambda_b' d\nu'(b')\) , we have, for every function \(h \in \mathcal{K}(T)\) ,

\[
\int h (t) d \mu (t) = \int d \nu^ {\prime} (b ^ {\prime}) \int h (t) d \lambda_ {b ^ {\prime}} ^ {\prime} (t) = \int d \nu (b) \int h (t) d \lambda_ {f (b)} ^ {\prime} (t)
\]

(Ch. V, §3, No.~4, Th. 1 and §6, No.~2, Th. 1); in other words, \(\mu = \int \lambda_{f(b)}' d\nu(b)\) . But since also \(\mu = \int \lambda_b d\nu(b)\) and since, for every \(b \in B - N_0\) , \(\lambda_b\) and \(\lambda_{f(b)}'\) are carried by \(\overline{p}^{-1}(b)\) , Th. 2 of No.~3 implies that \(\lambda_b = \lambda_{f(b)}'\) for almost every \(b \in B - N_0\) , hence for almost every \(b \in B\) . The conditions of Th. 4 are therefore verified by taking for \(N\) the union of \(N_0\) and the set of \(b \in B\) such that \(\lambda_b \neq \lambda_{f(b)}'\) .

\appendix
